\chapter{Existing AIOps Framework \& Problem Analysis}%
\label{ch:existing_framework}

% Doel: CloudSentinel (Jochem’s framework) beschrijven en analyseren als startpunt.

\section{Case context}
% Korte beschrijving van KPMG/bedrijf-context zoals in Jochem’s introductie.
This thesis is carried out in collaboration between the University of Amsterdam and Sue, a cloud provider and consultancy bureau based in Geldermalsen, Utrecht. At the university side, the work is supervised by Zhiming Zhao; on the company side, it is embedded in a managed-services context where teams already operate sizeable Kubernetes-based environments and have access to extensive monitoring data, but still spend substantial effort on manual root cause analysis and incident diagnosis. These operator pain points were discussed explicitly in earlier meetings and were a key motivation to move from pure observability tooling towards an AI-assisted approach for operations support.

The thesis builds directly on a previous Master’s project by Jochem van Gaalen, who designed and implemented CloudSentinel, an AIOps framework aimed at improving performance anomaly diagnosis in distributed cloud applications. CloudSentinel is positioned as a proof-of-concept framework that integrates machine-learning-based anomaly detection and root cause analysis into a Kubernetes and Prometheus-centred monitoring stack. That work provides the main technical starting point and reference implementation for the research presented in this thesis.

During the early phase of the project, the initial idea was to focus on abstractions and visualisation of service-level dependency graphs based on existing observability tools such as Istio/Jaeger and Cilium/Hubble. The goal there was to support operators in understanding complex service topologies and blast radius. After several discussions with the academic and industrial supervisors, this direction was reconsidered. The project was refocused towards AIOps and anomaly detection, with the explicit decision to use CloudSentinel as a baseline framework and to concentrate on automated monitoring, anomaly detection and root cause analysis rather than on pure topology visualisation.

Within this context, the core problem addressed by the present thesis can be summarised as follows: how well does an existing, research-oriented AIOps framework such as CloudSentinel transfer to a new microservice environment, and what are the practical and methodological limitations that still hinder its use for operator-facing anomaly diagnosis in real-world settings?

\section{Overview of the CloudSentinel-framework}
CloudSentinel is an AIOps framework designed to enhance performance monitoring, anomaly detection and root cause analysis for distributed cloud applications. The framework was developed as part of a Master’s thesis with the explicit aim of bridging the gap between theoretical advances in machine learning and their practical application in IT operations.It targets modern microservice architectures deployed on Kubernetes and assumes the presence of an existing monitoring stack, most notably Prometheus, which is used to collect time-series metrics from containers and services.

At a high level, CloudSentinel adopts a modular, microservice-oriented architecture. Different components are responsible for data collection, storage, processing, machine-learning analytics, and user interaction. A data collection module connects to monitoring backends such as Prometheus and ingests metric streams from the target environment. Configuration or mapping files are used to translate environment-specific metrics (e.g. Kubernetes pod CPU usage) into the internal feature representations required by the machine-learning models. In principle, these mappings allow the framework to be adapted to different deployments while keeping the core analytics components unchanged. The practical effort and limitations involved in this adaptation are evaluated later in this study.

Collected data is stored and processed in a way that supports both historical analysis and near real-time monitoring. A processing layer prepares the data for modelling by performing operations such as filtering, aggregation and normalisation. On top of this, the machine-learning layer hosts several algorithms. For anomaly detection, the framework integrates a CGNN-MHSA-AR model that operates on multivariate time-series data to detect deviations from learned normal behaviour. For root cause analysis, the framework incorporates a CausalRCA algorithm that aims to infer and rank causal relationships between metrics under anomalous conditions. Together, these components provide the core analytical capabilities of CloudSentinel.

The intended end-to-end CloudSentinel workflow can therefore be described as follows: Prometheus telemetry is collected from the target environment and prepared for analysis; the anomaly detection component processes the prepared data and produces anomaly indications; these indications can trigger the RCA component; and the resulting outputs are presented to the user through the web interface. This workflow, rather than the standalone performance of the embedded algorithms, forms the main unit of evaluation in this thesis.

The framework further includes a user-facing web interface implemented with standard web technologies. Through this interface, users can configure datasets and experiments, trigger training and evaluation runs, and inspect anomaly detection and RCA results. In experimental evaluation, CloudSentinel was applied to two main types of data: (1) a server-machine dataset (SMD) for anomaly detection experiments, and (2) a microservice-based demo application (Sock Shop) for evaluating the root cause analysis module. Performance was primarily assessed in terms of standard classification metrics such as precision, recall and F1-score for anomaly detection and ranking quality for RCA.

Finally, the accompanying thesis highlights the importance of continuous improvement and learning, and suggests further refinements of the embedded learning algorithms as well as extensions to other types of cloud environments and operational contexts as promising directions for future work.

\section{Known shortcomings of CloudSentinel}
% Dit gaat over de BEKENDE beperkingen uit Jochem’s thesi
% Dingen die hij zelf noemt: beperkte tests, geen volledig automatische response, beperkte dataset-diversiteit, etc.

The CloudSentinel thesis reports a number of limitations, both explicitly in discussions of results and future work and implicitly through the scope of the evaluation and user feedback. These limitations are important, because they frame where the present thesis can contribute.

% Generisability to other cloud environments, and and limited dataset use
First, the empirical evaluation of CloudSentinel is conducted on a limited set of environments and datasets. The anomaly detection experiments focus on the Server Machine Dataset (SMD), while root cause analysis is evaluated on a single microservice application (Sock Shop).Although the results demonstrate that the framework and its algorithms can operate effectively in these settings, the thesis itself acknowledges that further evaluation in other environments and under different operational conditions is needed. This raises open questions about the generalisability of the approach to other application domains, metric sets and workload characteristics.

% A lot of manual effort required, particularly in the training process
Second, the configuration and training processes in CloudSentinel still require a considerable amount of manual effort. The framework allows users to configure monitoring tasks and model parameters via its interface, and feedback from respondents indicates that the configuration process is generally perceived as usable. At the same time, user feedback points to the need for more automation and better support for advanced configuration options. In particular, suggestions include simplifying the training process by tighter integration with metrics backends such as Prometheus and automating aspects of data input and preparation. These comments imply that while the framework is functional, the data pipeline from raw metrics to training and evaluation datasets is not yet streamlined for everyday operator use.

% The training does not happen automatically; it needs automated training processes
Third, although the architecture is designed with continuous learning and adaptation in mind, the actual implementation of automated retraining cycles and user-driven anomaly labelling is still limited. The CloudSentinel thesis discusses strategies for integrating automated training processes and mechanisms for user labelling of anomalies as important directions for future work.The idea is that the framework should be able to evolve by incorporating new labelled incidents over time, improving detection accuracy and robustness. The fact that this is positioned as future work underscores that the current proof-of-concept still relies on largely offline, manually orchestrated training and evaluation runs.

% Enhancing the visualisation of cause-and-effect relationships in RCA
Fourth, there are usability and visualisation aspects that remain underdeveloped. The existing user study suggests that configuring monitoring tasks is relatively straightforward and that the overall interface is accessible. However, respondents also propose improvements such as offering more advanced configuration options, simplifying the training process, supporting the upload of custom models, and enhancing the visualisation of cause-and-effect relationships in RCA. These remarks indicate that, especially from an operator’s perspective, there is room to improve how anomaly detection and root cause analysis results are presented and how they fit into day-to-day workflows.

Finally, while the framework is intended to be scalable and robust, its evaluation is still largely limited to proof-of-concept scenarios. The thesis proposes further stress-testing under different configurations and loads, and extending experiments to more heterogeneous infrastructures as important next steps. This again highlights that the current version of CloudSentinel should be seen as a research prototype rather than a production-ready AIOps platform.

\section{Other observed limitations}
Beyond the limitations documented in the original thesis, a number of additional constraints and practical issues emerge when CloudSentinel is applied in the context of this project. These observations are based on attempts to deploy and use the framework on a new Kubernetes test environment running the Online Boutique microservice demo application, in combination with Prometheus and Grafana for metrics and dashboards.

% Over hoe SMD niet goed mapped op onze test applicatie. Roept de vraag op hoe generalizeerbaar CloudSentinel is
A first issue concerns portability and generalisation to a different application and metric set. CloudSentinel was originally tuned and evaluated on SMD and Sock Shop. When targeting Online Boutique, many of the default datasets, example configurations and mental models embedded in the framework no longer match the characteristics of the new environment. Metric naming, workload patterns and anomaly types all differ. As a result, substantial effort is required to define appropriate metric mappings and to construct training and testing datasets that are representative of “normal” and “anomalous” behaviour for Online Boutique. This reinforces the open question to what extent CloudSentinel can be used as a generic, environment-agnostic AIOps framework without significant re-engineering.

% De grote hoeveelheid manual werk bij het trainen van scratch, e.g. het moeten maken van custom scripts
A second practical limitation lies in the concrete data pipeline between Prometheus and CloudSentinel’s training and evaluation interfaces. The framework expects users to supply training arrays (typically normal-behaviour time series), testing arrays (time series including both normal and anomalous periods) and corresponding label arrays (ground truth per time step). In the current setup, constructing these arrays requires custom scripting to export metrics from Prometheus, align them in time, select relevant metrics and entities, and encode injected anomaly scenarios as labels. This design is powerful for controlled experiments but introduces friction and potential for error when onboarding a new environment, particularly compared to more integrated “click-to-train” experiences that operators might expect from commercial tools.

% Hoe het los/apart staat van andere dashboards zoals Grafana
Third, from a monitoring and user-experience perspective, CloudSentinel currently functions alongside existing dashboards rather than fully integrating with them. Low-level metrics and general system health are inspected in Grafana, whereas anomaly detection and RCA outputs are viewed in CloudSentinel’s own UI. The linkage between these views, e.g. from an anomaly alert to a rich, contextual view of affected services and metrics, is not yet seamless. For an operator whose workflow starts from a monitoring dashboard, this split can make it harder to use CloudSentinel as a natural extension of their existing tooling.

Together, these observations motivate the focus of the present thesis. The work aims to systematically reproduce and evaluate CloudSentinel in a different Kubernetes-based microservice environment, to quantify the configuration and data-preparation effort involved, and to analyse how well the framework supports anomaly diagnosis from an operator’s perspective. Based on this analysis, the thesis seeks to identify concrete opportunities to reduce manual effort, improve generalisability and strengthen the integration between anomaly detection, root cause analysis and operator-oriented visualisation.

\section{Other observed limitations}

Beyond the limitations documented in the original thesis, additional practical issues emerged while deploying and using CloudSentinel in a new Kubernetes environment. These observations were made during the initial setup of CloudSentinel for the Online Boutique microservice application, using Prometheus for metric collection and Grafana for monitoring dashboards. They provide the starting point for the evaluation in the following chapters.

% Over hoe SMD niet goed mapped op onze test applicatie. Roept de vraag op hoe generalizeerbaar CloudSentinel is
A first issue concerns portability to a different application and metric set. CloudSentinel was originally developed and evaluated using the SMD dataset and the Sock Shop application. When applying it to Online Boutique, the provided datasets, example configurations, and metric assumptions did not directly match the new environment. Metric names, available services, workload behaviour, and expected failure signals differed. As a result, environment-specific metric mappings and dataset definitions had to be created manually. This raises the question of how easily CloudSentinel can be adopted in a new environment without substantial re-engineering.

% De grote hoeveelheid manual werk bij het trainen van scratch, e.g. het moeten maken van custom scripts
A second issue concerns the data preparation pipeline. CloudSentinel expects training data, test data, and time-step labels in specific array formats. In the current setup, these inputs must be created through custom scripts that export metrics from Prometheus, align time series, select relevant entities, and encode known fault periods as labels. This process supports controlled experimentation, but it also requires considerable manual effort and creates opportunities for inconsistencies between training, offline evaluation, and live monitoring.

% Het moeten handmatig slepen van data naar pods
A third issue concerns configuration and reproducibility. Some configuration changes made through the user interface did not automatically propagate to the components that perform monitoring or inference. In addition, trained model artefacts and environment-specific settings had to be transferred or updated manually. These behaviours make the deployment process harder to reproduce and suggest that important configuration is spread across multiple locations.

% Hoe het los/apart staat van andere dashboards zoals Grafana
A fourth issue concerns the operator workflow. Grafana is used to inspect low-level metrics and general system health, while CloudSentinel provides anomaly detection and RCA outputs through a separate interface. The connection between these views is limited. An operator must therefore manually relate a detected anomaly or RCA result to the relevant services, metrics, and historical context in Grafana.

% De limitaties van hoe RCA wordt gepresenteerd
Finally, the current RCA output is primarily technical and metric-oriented. Although it can provide a ranking of potentially relevant variables, the output does not always clearly identify the affected service or explain the result in a form that is directly useful to an operator. This limits the practical interpretability of the diagnosis workflow.

These observations define the focus of the baseline evaluation. Chapters~\ref{ch:requirements_for_framework}--\ref{ch:baseline_evaluation_results} assess whether CloudSentinel supports the end-to-end anomaly detection and RCA workflow, how much effort is required to adapt it to a new environment, how reproducible its configuration is, and how usable its outputs are for an operator. The findings are then used to derive requirements for a focused framework improvement.